\nonstopmode{}
\documentclass[a4paper]{book}
\usepackage[times,inconsolata,hyper]{Rd}
\usepackage{makeidx}
\makeatletter\@ifl@t@r\fmtversion{2018/04/01}{}{\usepackage[utf8]{inputenc}}\makeatother
% \usepackage{graphicx} % @USE GRAPHICX@
\makeindex{}
\begin{document}
\chapter*{}
\begin{center}
{\textbf{\huge Package `collaborativedatasciencepackage'}}
\par\bigskip{\large \today}
\end{center}
\ifthenelse{\boolean{Rd@use@hyper}}{\hypersetup{pdftitle = {collaborativedatasciencepackage: Analysis Tools for Coral Reef Benthic Cover Data}}}{}
\ifthenelse{\boolean{Rd@use@hyper}}{\hypersetup{pdfauthor = {Juan Pablo Lozano Pena; Megan Do; Carly Dempsey}}}{}
\begin{description}
\raggedright{}
\item[Title]\AsIs{Analysis Tools for Coral Reef Benthic Cover Data}
\item[Version]\AsIs{0.0.0.9000}
\item[Description]\AsIs{Provides data, analysis workflows, and visualization tools
for examining benthic community structure in coral reef ecosystems.
Supports functional group aggregation, ordination using Principal
Coordinate Analysis (PCoA), and Bayesian generalized linear modelling
to compare reef habitats and coral--macroalgae relationships.}
\item[License]\AsIs{MIT + file LICENSE}
\item[Encoding]\AsIs{UTF-8}
\item[Depends]\AsIs{R (>= 4.1.0)}
\item[Imports]\AsIs{BiodiversityR, brms, dplyr, ggdist, ggplot2, ggrepel, modelr,
patchwork, tibble, tidybayes, vegan}
\item[Suggests]\AsIs{knitr, rmarkdown, tidyverse}
\item[VignetteBuilder]\AsIs{knitr}
\item[Roxygen]\AsIs{list(markdown = TRUE)}
\item[RoxygenNote]\AsIs{7.3.3}
\item[LazyData]\AsIs{true}
\item[NeedsCompilation]\AsIs{no}
\item[Author]\AsIs{Juan Pablo Lozano Pena [aut],
Megan Do [aut],
Carly Dempsey [aut, cre]}
\item[Maintainer]\AsIs{Carly Dempsey }\email{cd39227@my.utexas.edu}\AsIs{}
\end{description}
\Rdcontents{Contents}
\HeaderA{collaborativedatasciencepackage-package}{collaborativedatascience-package}{collaborativedatasciencepackage.Rdash.package}
\aliasA{collaborativedatasciencepackage}{collaborativedatasciencepackage-package}{collaborativedatasciencepackage}
\keyword{internal}{collaborativedatasciencepackage-package}
%
\begin{Description}
A package for analyzing benthic cover data from coral reef ecosystems,
including data wrangling, visualization, ordination, and modeling tools
developed for a collaborative data science project.
\end{Description}
%
\begin{Details}
The collaborativedatasciencepackage provides functions and datasets
to support ecological analysis of coral reef benthic communities.
Key features include:

\begin{itemize}

\item{} Data preparation and wrangling of quadrat-level observations
\item{} Functional group aggregation of benthic categories
\item{} Multivariate ordination (PCoA)
\item{} Visualization utilities for benthic community data
\item{} Modeling tools for coral–macroalgae relationships

\end{itemize}

\end{Details}
%
\begin{Section}{Package options}

This package does not currently define any global options.
\end{Section}
%
\begin{Author}
Carly R. Dempsey
\end{Author}
\HeaderA{benthic\_plot}{Plot of Benthic PCoA Analysis}{benthic.Rul.plot}
%
\begin{Description}
This function plots the PCoA run previously, adding elements that visualize factors habitat coordinates and species.
\end{Description}
%
\begin{Usage}
\begin{verbatim}
benthic_plot(benthic.species.scores, benthic.site.scores, benthic.hulls)
\end{verbatim}
\end{Usage}
%
\begin{Arguments}
\begin{ldescription}
\item[\code{benthic.species.scores}] A data frame of species scores

\item[\code{benthic.site.scores}] A data frame of site scores

\item[\code{benthic.hulls}] A data frame containing convex hull coordinates
\end{ldescription}
\end{Arguments}
%
\begin{Value}
A ggplot object
\end{Value}
\HeaderA{habitat\_effect}{Habitat Effect on Coral Cover}{habitat.Rul.effect}
%
\begin{Description}
This function determines the interaction between coral cover and habitat of our Bayesian GLM
\end{Description}
%
\begin{Usage}
\begin{verbatim}
habitat_effect(data, model)
\end{verbatim}
\end{Usage}
%
\begin{Arguments}
\begin{ldescription}
\item[\code{data}] data object being used

\item[\code{model}] fitted Bayesian model object being used
\end{ldescription}
\end{Arguments}
%
\begin{Value}
A ggplot object
\end{Value}
\printindex{}
\end{document}
